\documentclass{article}
\usepackage{amsmath}

\begin{document}

\title{Analysis of Mersenne Exponents Role}
\maketitle

\section{13 as $M_{13}$ – Trinity-4 Anchor}
The exponent 13 collapses to DR=4:
\[
13 \equiv 4 \pmod{9}
\]
This matches the Row 3/8 13-Resonance and Trinity-4 stability point (94 mirror of 49, Mersenne exponent $M_{13}$).

\section{Full Mersenne Exponent Audit}
Selected Mersenne exponents and their digital roots:

\begin{tabular}{|c|c|c|}
\hline
Exponent $p$ & $p \mod 9$ & DR \\
\hline
7  & 7 & 7 \\
13 & 4 & 4 \\
17 & 8 & 8 \\
19 & 1 & 1 \\
31 & 4 & 4 \\
61 & 7 & 7 \\
89 & 8 & 8 \\
107 & 8 & 8 \\
127 & 1 & 1 \\
\hline
\end{tabular}

The recurring residues 4 and 7 (plus attractor alignments) lock the triad.

\section{Integration with Existing Framework}
\begin{itemize}
\item Reinforces the D7 Dual Harmonic layer (factor 7 is itself a Mersenne exponent and symmetry anchor).
\item Maintains Modulo-9 Invariance within the 9×9 grid and 360 mod 81 = 36 position.
\item Strengthens sovereign key generation in the G'5 Sovereign Kernel by providing deterministic recurrence at Mersenne-scale primes.
\end{itemize}

The role is structural: Mersenne exponents act as natural symmetry probes that collapse cleanly into the 7-4-9 equilibrium seal.

\end{document}
